\documentclass[12pt]{article}
\usepackage[T2A]{fontenc}
\usepackage[utf8]{inputenc}
\usepackage[russian]{babel}
\usepackage{amsmath,amssymb}
\usepackage{geometry}
\usepackage{graphicx}
\usepackage{booktabs}
\usepackage{enumitem}
\usepackage{caption}
\usepackage{tikz}
\usetikzlibrary{positioning,shapes,arrows.meta,fit,backgrounds,calc}
\geometry{a4paper, margin=2.3cm}
\captionsetup{font=small,labelfont=bf}

\title{Пайплайн генерации обучающих данных\\
для нейрооператоров прогнозирования\\
термоупругих волн в геологических средах}
\author{}
\date{}

\begin{document}
\maketitle

\section*{Цель}

Построить единый «мастер-датасет», из которого без перегенерации можно
обучать четыре класса моделей: \textbf{PINN}, \textbf{FNO},
\textbf{MeshGraphNet} и \textbf{Transformer-нейрооператор}. Целевая задача~---
аппроксимация оператора решения связанной системы термоупругости
\[
\mathcal{G}: \bigl(a(x),\; \mathrm{IC},\; \mathrm{BC},\; s(x,t)\bigr)
\;\longmapsto\;
\bigl(u(x,t),\; T(x,t),\; \sigma(x,t)\bigr),
\]
где $a(x)=(\rho,\lambda,\mu,\gamma,k,c_p)$~--- поля материальных параметров
неоднородной (слоистой) среды, $s$~--- источник, $\mathrm{IC},\mathrm{BC}$~---
начальные и граничные условия.

\section{Источник данных: COMSOL Multiphysics}

В качестве «решателя истины» используется COMSOL~5.4 со связкой модулей
\textit{Solid Mechanics} и \textit{Heat Transfer in Solids} с термоупругой
связкой (\textit{Thermal Expansion}). Такая постановка соответствует системе
Куна\,--\,Томпсона, использованной в работе научного руководителя:
\begin{align}
\rho \frac{\partial^2 u_i}{\partial t^2} &= \frac{\partial \sigma_{ij}}{\partial x_j},\\
\sigma_{ij} &= \lambda\,\delta_{ij}\varepsilon_{kk} + 2\mu\,\varepsilon_{ij}
              - \gamma\,\delta_{ij}\,T,\\
k\,\nabla^2 T - C\,\frac{\partial T}{\partial t}
              &= \gamma\,T_0\,\frac{\partial \varepsilon_{kk}}{\partial t}.
\end{align}
COMSOL даёт численно консистентные поля $u, T, \sigma_{ij}, \varepsilon_{ij}$
на неструктурированном тетраэдральном меше, что используется как
ground truth для обучения нейрооператоров.

\section{Параметризация задачи}

В файле модели \texttt{Untitled.mph} вынесены в \emph{Global Parameters}
все характеристики, которые варьируются от сэмпла к сэмплу:
\begin{itemize}[leftmargin=1.4em,itemsep=2pt]
\item \textbf{Среда:} число слоёв $L\in\{2,\dots,6\}$, толщины слоёв,
      материал каждого слоя (выбирается из таблицы пород,
      см. \texttt{Insar/data/}).
\item \textbf{Источник:} тип (механический импульс, тепловой импульс,
      комбинированный), позиция $x_s$, амплитуда $A$, частотный профиль
      (Ricker wavelet с центральной частотой $f_0$).
\item \textbf{Граничные условия:} тип на каждой грани области
      (Dirichlet, Neumann, поглощающие).
\end{itemize}
Случайный генератор конфигураций \texttt{sample\_random\_medium.py}
производит $N$ независимых конфигураций, которые передаются в
\textit{Parametric Sweep} COMSOL.

\begin{table}[h]
\centering
\small
\begin{tabular}{lll}
\toprule
\textbf{Параметр} & \textbf{Распределение} & \textbf{Источник значений}\\
\midrule
число слоёв $L$ & $\mathrm{Uniform}\{2,\dots,6\}$ & ---\\
толщины $h_\ell$ & Дирихле, $\sum h_\ell = H$ & ---\\
материал слоя & категориально & \texttt{materials.csv}\\
$\rho,\lambda,\mu,\gamma,k,c_p$ & по материалу & таблицы пород Insar\\
позиция источника $x_s$ & $\mathrm{Uniform}$ в области & ---\\
амплитуда $A$, частота $f_0$ & log-uniform & ---\\
тип источника & $\{$mech,\,thermal,\,coupled$\}$ & ---\\
\bottomrule
\end{tabular}
\caption{Случайные параметры, варьируемые между сэмплами датасета.}
\end{table}

\section{Структура одного сэмпла}

Каждый узел расчётной сетки описывается \textbf{24 каналами},
сгруппированными по физическому смыслу
(см.~\texttt{feature\_names.npy} текущего гранитного сэмпла):

\begin{figure}[h]
\centering
% Структура одного сэмпла: 24 канала на узел, сгруппированные по физическому смыслу
\begin{tikzpicture}[
  font=\small,
  node distance=2mm and 4mm,
  ch/.style={
    rectangle, rounded corners=1.5pt,
    draw=black!55, fill=white,
    minimum width=18mm, minimum height=7mm,
    inner sep=2pt, font=\scriptsize
  },
  grp/.style={
    rectangle, rounded corners=3pt,
    inner sep=3mm, draw=black!40, dashed
  },
  hdr/.style={font=\bfseries\small}
]

% Group: thermal
\node[ch, fill=red!10]   (T)  {$T$};
\node[ch, fill=red!10, right=of T] (k)  {$k$};
\node[ch, fill=red!10, right=of k] (rh) {$\rho_h$};
\node[ch, fill=red!10, right=of rh] (Cp) {$C_p$};

\node[grp, fit=(T)(k)(rh)(Cp), label={[hdr]above:Тепловые поля и параметры (4)}] (gT) {};

% Group: displacement + velocity
\node[ch, fill=blue!10, below=14mm of T]    (u)  {$u$};
\node[ch, fill=blue!10, right=of u]         (v)  {$v$};
\node[ch, fill=blue!10, right=of v]         (w)  {$w$};
\node[ch, fill=blue!10, right=of w]         (ut) {$u_t$};
\node[ch, fill=blue!10, right=of ut]        (vt) {$v_t$};
\node[ch, fill=blue!10, right=of vt]        (wt) {$w_t$};

\node[grp, fit=(u)(v)(w)(ut)(vt)(wt), label={[hdr]above:Смещение и скорость (6)}] (gU) {};

% Group: stress
\node[ch, fill=violet!12, below=14mm of u] (mises) {mises};
\node[ch, fill=violet!12, right=of mises]  (sx) {$\sigma_x$};
\node[ch, fill=violet!12, right=of sx]     (sy) {$\sigma_y$};
\node[ch, fill=violet!12, right=of sy]     (sz) {$\sigma_z$};
\node[ch, fill=violet!12, right=of sz]     (sxy){$\sigma_{xy}$};
\node[ch, fill=violet!12, right=of sxy]    (syz){$\sigma_{yz}$};
\node[ch, fill=violet!12, right=of syz]    (sxz){$\sigma_{xz}$};

\node[grp, fit=(mises)(sx)(sy)(sz)(sxy)(syz)(sxz), label={[hdr]above:Напряжения (7)}] (gS) {};

% Group: strain
\node[ch, fill=teal!12, below=14mm of mises] (eX) {$\varepsilon_X$};
\node[ch, fill=teal!12, right=of eX]         (eY) {$\varepsilon_Y$};
\node[ch, fill=teal!12, right=of eY]         (eZ) {$\varepsilon_Z$};

\node[grp, fit=(eX)(eY)(eZ), label={[hdr]above:Деформации диагональные (3)}] (gE) {};

% Group: material params
\node[ch, fill=orange!15, right=24mm of eZ] (E)   {$E$};
\node[ch, fill=orange!15, right=of E]       (nu)  {$\nu$};
\node[ch, fill=orange!15, right=of nu]      (rs)  {$\rho_s$};
\node[ch, fill=orange!15, right=of rs]      (al)  {$\alpha$};

\node[grp, fit=(E)(nu)(rs)(al), label={[hdr]above:Материальные параметры (4)}] (gM) {};

\end{tikzpicture}

\caption{24 канала состояния узла, сгруппированные по физическому смыслу:
тепловые (4), кинематические (6), напряжения (7), деформации (3),
материальные параметры (4).}
\label{fig:sample}
\end{figure}

\noindent
Эти 24 канала образуют \emph{избыточный} набор: для PDE-residual в PINN
и для предсказания достаточно подмножества, остальное служит
вспомогательной диагностикой и регуляризацией.

\section{Препроцессинг: от COMSOL CSV к графовому представлению}

Сырые экспортируемые из COMSOL файлы имеют ширину
$3 + n_{\mathrm{ch}}\cdot N_t$ столбцов
(координаты узла $+$ канал на каждом временном шаге).
Скрипт \texttt{convert\_raw\_to\_graph.py} выполняет:
\begin{enumerate}[leftmargin=1.4em,itemsep=2pt]
\item чтение \texttt{mesh.mphtxt} $\rightarrow$ массив координат
      \texttt{coords} формы $(N,3)$;
\item построение рёберного списка \texttt{edge\_index} формы $(E,2)$
      по тетраэдральным симплексам, удаление дубликатов;
\item вычисление признаков рёбер
      $(\Delta x, \Delta y, \Delta z, \|\Delta\|)$
      $\rightarrow$ \texttt{edge\_attr} формы $(E,4)$;
\item слияние всех CSV по идентификатору узла и времени
      $\rightarrow$ \texttt{node\_features} формы $(T,N,24)$;
\item вычисление per-channel статистик
      \texttt{norm\_mean}, \texttt{norm\_std}
      $\rightarrow$ нормированный тензор \texttt{node\_features}
      и сырой \texttt{node\_features\_raw}.
\end{enumerate}

В текущей реализации \emph{пилотного} гранитного сэмпла получены тензоры
\texttt{coords}~$(1020,3)$,
\texttt{edge\_index}~$(11422,2)$,
\texttt{edge\_attr}~$(11422,4)$,
\texttt{node\_features}~$(3001,1020,24)$,
\texttt{times}~$(3001,)$,
с шагом по времени $\Delta t = 0.01$~с на интервале $t\in[0,30]$~с.

\section{Мастер-датасет: структура папок}

Сводный датасет имеет следующую структуру:
\begin{verbatim}
thermoelastic_3d_v1/
  samples/
    sample_00000/
      coords.npy           # (N, 3)
      edge_index.npy       # (E, 2)
      edge_attr.npy        # (E, 4)
      node_features.npy    # (T, N, 24), нормированные
      node_features_raw.npy
      times.npy            # (T,)
      norm_mean.npy        # (24,)
      norm_std.npy         # (24,)
      medium.json          # описание слоёв
      source.json          # источник
      bc.json              # граничные условия
    sample_00001/
    ...
  materials.csv            # таблица пород (из Insar/data)
  meta.json                # уравнения, единицы, диапазоны
  splits.json              # train / val / test
\end{verbatim}

Раздельные файлы \texttt{medium.json}, \texttt{source.json},
\texttt{bc.json} нужны, чтобы модели могли использовать конфигурацию
сэмпла как условие (conditioning), а не «угадывать» источник из полей.

\section{Адаптация под четыре модели}

Из одного и того же сэмпла строятся четыре различных представления
с помощью даталоадеров.

\begin{figure}[h]
\centering
% Как один сэмпл превращается в данные для каждой из четырёх архитектур
\begin{tikzpicture}[
  font=\small,
  node distance=4mm and 8mm,
  src/.style={
    rectangle, rounded corners=2pt,
    draw=violet!60!black, fill=violet!10,
    minimum width=64mm, minimum height=12mm,
    inner sep=4pt, align=center
  },
  view/.style={
    rectangle, rounded corners=2pt,
    draw=teal!60!black, fill=teal!8,
    minimum width=42mm, minimum height=24mm,
    inner sep=4pt, align=center
  },
  ttl/.style={font=\bfseries\small},
  arr/.style={-{Latex[length=2mm]}, thick, draw=black!60}
]

\node[src] (s) {Сэмпл (мастер): \texttt{node\_features} (T,N,24) + \texttt{coords} + \texttt{edges} + \texttt{times}};

% PINN
\node[view, below left=12mm and 32mm of s] (pinn)
  {\ttl PINN\\[2pt]
   \scriptsize случайные точки\\$\{(x_i, t_i)\}$:\\
   collocation, BC, IC,\\
   data-loss};

% FNO
\node[view, below left=12mm and -2mm of s] (fno)
  {\ttl FNO\\[2pt]
   \scriptsize интерполяция\\на регулярную\\сетку\\$(N_x, N_y, N_z, N_t)$};

% MGN
\node[view, below right=12mm and -2mm of s] (mgn)
  {\ttl MeshGraphNet\\[2pt]
   \scriptsize узлы + рёбра,\\окно $t \to t{+}1$,\\автoрегрессионно};

% Transformer
\node[view, below right=12mm and 32mm of s] (tr)
  {\ttl Transformer\\[2pt]
   \scriptsize input-токены\\$(p_i, a_i, \mathrm{IC}_i)$,\\query-токены\\$(x_q, t_q)$};

\draw[arr] (s.south) -- ($(pinn.north)+(0,2mm)$);
\draw[arr] (s.south) -- ($(fno.north)+(0,2mm)$);
\draw[arr] (s.south) -- ($(mgn.north)+(0,2mm)$);
\draw[arr] (s.south) -- ($(tr.north)+(0,2mm)$);

\end{tikzpicture}

\caption{Преобразование одного мастер-сэмпла в формат, ожидаемый каждой из
четырёх архитектур.}
\label{fig:loaders}
\end{figure}

\paragraph{PINN-loader.} Возвращает три (или четыре) набора точек:
$N_{\mathrm{col}}$ случайных точек $(x_i, t_i)$ внутри области
для PDE-residual; $N_{\mathrm{bc}}$ точек на границе с истинными значениями;
$N_{\mathrm{ic}}$ точек при $t=0$; опционально $N_{\mathrm{data}}$ точек
из \texttt{node\_features} как data-loss. Параметры среды интерполируются
с \texttt{coords} в эти точки.

\paragraph{FNO-loader.} Интерполирует \texttt{node\_features} с
неструктурированного меша на регулярную сетку
$(N_x, N_y, N_z)$~--- это требование FFT в спектральных слоях FNO.
Вход: тензор $(C_{\mathrm{in}}, N_x, N_y, N_z)$ из материальных параметров
и начальных условий; выход: $(C_{\mathrm{out}}, N_x, N_y, N_z, N_t)$
из полей $u, T, \sigma$.

\paragraph{MGN-loader.} Возвращает PyTorch Geometric \texttt{Data}-объект:
$\mathtt{x} = $ \texttt{node\_features}$[t]$,
$\mathtt{edge\_index} = $ \texttt{edge\_index.T},
$\mathtt{edge\_attr} = $ \texttt{edge\_attr},
$\mathtt{y} = $ \texttt{node\_features}$[t{+}1] - $\texttt{node\_features}$[t]$
(приращение, более устойчиво).

\paragraph{Transformer-loader.} Случайно сэмплирует
$N_{\mathrm{in}}$~входных токенов и $N_{\mathrm{q}}$~запросных токенов.
Признаки входного токена~--- $(p_i, a_i, \mathrm{IC}_i)$, запрос~---
координата $(x_q, t_q)$, целевое значение~--- поля $(u, T, \sigma)$ в
этой координате. Разрешение свободное: можно тренировать на 256 точках
и инферить на 4096.

\section{Полный воркфлоу}

\begin{figure}[h]
\centering
% Главная диаграмма пайплайна генерации данных
% Подключается в main.tex через % Главная диаграмма пайплайна генерации данных
% Подключается в main.tex через % Главная диаграмма пайплайна генерации данных
% Подключается в main.tex через \input{fig_pipeline}
\begin{tikzpicture}[
  font=\small,
  node distance=6mm and 10mm,
  every node/.style={align=center},
  input/.style={
    rectangle, rounded corners=2pt,
    draw=orange!70!black, fill=orange!10,
    minimum width=42mm, minimum height=10mm,
    inner sep=4pt
  },
  proc/.style={
    rectangle, rounded corners=2pt,
    draw=blue!60!black, fill=blue!8,
    minimum width=72mm, minimum height=11mm,
    inner sep=4pt
  },
  data/.style={
    rectangle, rounded corners=2pt,
    draw=black!50, fill=black!5,
    minimum width=72mm, minimum height=12mm,
    inner sep=4pt
  },
  master/.style={
    rectangle, rounded corners=2pt,
    draw=violet!60!black, fill=violet!10,
    minimum width=72mm, minimum height=11mm,
    inner sep=4pt
  },
  loader/.style={
    rectangle, rounded corners=2pt,
    draw=teal!60!black, fill=teal!10,
    minimum width=28mm, minimum height=10mm,
    inner sep=3pt
  },
  model/.style={
    rectangle, rounded corners=2pt,
    draw=green!50!black, fill=green!12,
    minimum width=28mm, minimum height=9mm,
    inner sep=3pt, font=\small\bfseries
  },
  arr/.style={-{Latex[length=2mm,width=2mm]}, thick, draw=black!65},
  arrthin/.style={-{Latex[length=1.6mm,width=1.6mm]}, draw=black!55}
]

% --- Stage 1: inputs ---
\node[input] (mat)   {Таблица свойств пород\\\scriptsize{(Insar/data, гранит/базальт/...)}};
\node[input, right=of mat] (rules) {Правила сэмплирования\\\scriptsize{слои, источник, ГУ}};

% --- Stage 2: per-sample config ---
\node[proc, below=8mm of $(mat)!0.5!(rules)$] (cfg)
  {Конфигурация сэмпла $i$\\\scriptsize{$\{$слои, материалы, источник, ГУ$\}_i$}};

\draw[arr] (mat) -- (cfg);
\draw[arr] (rules) -- (cfg);

% --- Stage 3: COMSOL ---
\node[proc, below=of cfg] (comsol)
  {COMSOL \texttt{Untitled.mph} + Global Parameters\\
   \scriptsize{Parametric Sweep $\rightarrow$ N независимых решений}};
\draw[arr] (cfg) -- (comsol);

% --- Stage 4: raw export ---
\node[data, below=of comsol] (raw)
  {\textbf{Raw export на сэмпл} (\texttt{raw data/}):\\
   \scriptsize{\texttt{data\_displacement.csv},
   \texttt{data\_temperature.csv},
   \texttt{data\_stress\_\{1,2,3\}.csv},}\\
   \scriptsize{\texttt{data\_materials.csv}, \texttt{mesh.mphtxt}}};
\draw[arr] (comsol) -- (raw);

% --- Stage 5: preprocessing ---
\node[proc, below=of raw] (prep)
  {Препроцессинг \texttt{convert\_raw\_to\_graph.py}\\
   \scriptsize{разбор CSV $\to$ узлы, рёбра, признаки, нормализация}};
\draw[arr] (raw) -- (prep);

% --- Stage 6: per-sample npy ---
\node[data, below=of prep] (npy)
  {\textbf{Графовое представление сэмпла} (\texttt{data not raw/}):\\
   \scriptsize{\texttt{coords.npy} (N,3),\;
   \texttt{edge\_index.npy} (E,2),\;
   \texttt{edge\_attr.npy} (E,4),}\\
   \scriptsize{\texttt{node\_features.npy} (T,N,24),\;
   \texttt{times.npy} (T,),\;
   \texttt{norm\_mean/std.npy}}};
\draw[arr] (prep) -- (npy);

% --- Stage 7: master ---
\node[master, below=of npy] (master)
  {\textbf{Мастер-датасет}\\
   \scriptsize{\texttt{samples/sample\_XXXXX/} $\times$ N,\;
   \texttt{materials.csv},\; \texttt{meta.json},\; \texttt{splits.json}}};
\draw[arr] (npy) -- (master);

% --- Stage 8: four loaders ---
\node[loader, below=12mm of master, xshift=-48mm] (l1) {PINN-loader\\\scriptsize{collocation/BC/IC}};
\node[loader, below=12mm of master, xshift=-16mm] (l2) {FNO-loader\\\scriptsize{регулярная сетка}};
\node[loader, below=12mm of master, xshift=16mm]  (l3) {MGN-loader\\\scriptsize{узлы + рёбра, $t \to t{+}1$}};
\node[loader, below=12mm of master, xshift=48mm]  (l4) {Transformer-loader\\\scriptsize{input/query токены}};

\draw[arr] (master.south) -- ($(l1.north)+(0,2mm)$);
\draw[arr] (master.south) -- ($(l2.north)+(0,2mm)$);
\draw[arr] (master.south) -- ($(l3.north)+(0,2mm)$);
\draw[arr] (master.south) -- ($(l4.north)+(0,2mm)$);

% --- Stage 9: models ---
\node[model, below=6mm of l1] (m1) {PINN};
\node[model, below=6mm of l2] (m2) {FNO};
\node[model, below=6mm of l3] (m3) {MeshGraphNet};
\node[model, below=6mm of l4] (m4) {Transformer};

\draw[arrthin] (l1) -- (m1);
\draw[arrthin] (l2) -- (m2);
\draw[arrthin] (l3) -- (m3);
\draw[arrthin] (l4) -- (m4);

\end{tikzpicture}

\begin{tikzpicture}[
  font=\small,
  node distance=6mm and 10mm,
  every node/.style={align=center},
  input/.style={
    rectangle, rounded corners=2pt,
    draw=orange!70!black, fill=orange!10,
    minimum width=42mm, minimum height=10mm,
    inner sep=4pt
  },
  proc/.style={
    rectangle, rounded corners=2pt,
    draw=blue!60!black, fill=blue!8,
    minimum width=72mm, minimum height=11mm,
    inner sep=4pt
  },
  data/.style={
    rectangle, rounded corners=2pt,
    draw=black!50, fill=black!5,
    minimum width=72mm, minimum height=12mm,
    inner sep=4pt
  },
  master/.style={
    rectangle, rounded corners=2pt,
    draw=violet!60!black, fill=violet!10,
    minimum width=72mm, minimum height=11mm,
    inner sep=4pt
  },
  loader/.style={
    rectangle, rounded corners=2pt,
    draw=teal!60!black, fill=teal!10,
    minimum width=28mm, minimum height=10mm,
    inner sep=3pt
  },
  model/.style={
    rectangle, rounded corners=2pt,
    draw=green!50!black, fill=green!12,
    minimum width=28mm, minimum height=9mm,
    inner sep=3pt, font=\small\bfseries
  },
  arr/.style={-{Latex[length=2mm,width=2mm]}, thick, draw=black!65},
  arrthin/.style={-{Latex[length=1.6mm,width=1.6mm]}, draw=black!55}
]

% --- Stage 1: inputs ---
\node[input] (mat)   {Таблица свойств пород\\\scriptsize{(Insar/data, гранит/базальт/...)}};
\node[input, right=of mat] (rules) {Правила сэмплирования\\\scriptsize{слои, источник, ГУ}};

% --- Stage 2: per-sample config ---
\node[proc, below=8mm of $(mat)!0.5!(rules)$] (cfg)
  {Конфигурация сэмпла $i$\\\scriptsize{$\{$слои, материалы, источник, ГУ$\}_i$}};

\draw[arr] (mat) -- (cfg);
\draw[arr] (rules) -- (cfg);

% --- Stage 3: COMSOL ---
\node[proc, below=of cfg] (comsol)
  {COMSOL \texttt{Untitled.mph} + Global Parameters\\
   \scriptsize{Parametric Sweep $\rightarrow$ N независимых решений}};
\draw[arr] (cfg) -- (comsol);

% --- Stage 4: raw export ---
\node[data, below=of comsol] (raw)
  {\textbf{Raw export на сэмпл} (\texttt{raw data/}):\\
   \scriptsize{\texttt{data\_displacement.csv},
   \texttt{data\_temperature.csv},
   \texttt{data\_stress\_\{1,2,3\}.csv},}\\
   \scriptsize{\texttt{data\_materials.csv}, \texttt{mesh.mphtxt}}};
\draw[arr] (comsol) -- (raw);

% --- Stage 5: preprocessing ---
\node[proc, below=of raw] (prep)
  {Препроцессинг \texttt{convert\_raw\_to\_graph.py}\\
   \scriptsize{разбор CSV $\to$ узлы, рёбра, признаки, нормализация}};
\draw[arr] (raw) -- (prep);

% --- Stage 6: per-sample npy ---
\node[data, below=of prep] (npy)
  {\textbf{Графовое представление сэмпла} (\texttt{data not raw/}):\\
   \scriptsize{\texttt{coords.npy} (N,3),\;
   \texttt{edge\_index.npy} (E,2),\;
   \texttt{edge\_attr.npy} (E,4),}\\
   \scriptsize{\texttt{node\_features.npy} (T,N,24),\;
   \texttt{times.npy} (T,),\;
   \texttt{norm\_mean/std.npy}}};
\draw[arr] (prep) -- (npy);

% --- Stage 7: master ---
\node[master, below=of npy] (master)
  {\textbf{Мастер-датасет}\\
   \scriptsize{\texttt{samples/sample\_XXXXX/} $\times$ N,\;
   \texttt{materials.csv},\; \texttt{meta.json},\; \texttt{splits.json}}};
\draw[arr] (npy) -- (master);

% --- Stage 8: four loaders ---
\node[loader, below=12mm of master, xshift=-48mm] (l1) {PINN-loader\\\scriptsize{collocation/BC/IC}};
\node[loader, below=12mm of master, xshift=-16mm] (l2) {FNO-loader\\\scriptsize{регулярная сетка}};
\node[loader, below=12mm of master, xshift=16mm]  (l3) {MGN-loader\\\scriptsize{узлы + рёбра, $t \to t{+}1$}};
\node[loader, below=12mm of master, xshift=48mm]  (l4) {Transformer-loader\\\scriptsize{input/query токены}};

\draw[arr] (master.south) -- ($(l1.north)+(0,2mm)$);
\draw[arr] (master.south) -- ($(l2.north)+(0,2mm)$);
\draw[arr] (master.south) -- ($(l3.north)+(0,2mm)$);
\draw[arr] (master.south) -- ($(l4.north)+(0,2mm)$);

% --- Stage 9: models ---
\node[model, below=6mm of l1] (m1) {PINN};
\node[model, below=6mm of l2] (m2) {FNO};
\node[model, below=6mm of l3] (m3) {MeshGraphNet};
\node[model, below=6mm of l4] (m4) {Transformer};

\draw[arrthin] (l1) -- (m1);
\draw[arrthin] (l2) -- (m2);
\draw[arrthin] (l3) -- (m3);
\draw[arrthin] (l4) -- (m4);

\end{tikzpicture}

\begin{tikzpicture}[
  font=\small,
  node distance=6mm and 10mm,
  every node/.style={align=center},
  input/.style={
    rectangle, rounded corners=2pt,
    draw=orange!70!black, fill=orange!10,
    minimum width=42mm, minimum height=10mm,
    inner sep=4pt
  },
  proc/.style={
    rectangle, rounded corners=2pt,
    draw=blue!60!black, fill=blue!8,
    minimum width=72mm, minimum height=11mm,
    inner sep=4pt
  },
  data/.style={
    rectangle, rounded corners=2pt,
    draw=black!50, fill=black!5,
    minimum width=72mm, minimum height=12mm,
    inner sep=4pt
  },
  master/.style={
    rectangle, rounded corners=2pt,
    draw=violet!60!black, fill=violet!10,
    minimum width=72mm, minimum height=11mm,
    inner sep=4pt
  },
  loader/.style={
    rectangle, rounded corners=2pt,
    draw=teal!60!black, fill=teal!10,
    minimum width=28mm, minimum height=10mm,
    inner sep=3pt
  },
  model/.style={
    rectangle, rounded corners=2pt,
    draw=green!50!black, fill=green!12,
    minimum width=28mm, minimum height=9mm,
    inner sep=3pt, font=\small\bfseries
  },
  arr/.style={-{Latex[length=2mm,width=2mm]}, thick, draw=black!65},
  arrthin/.style={-{Latex[length=1.6mm,width=1.6mm]}, draw=black!55}
]

% --- Stage 1: inputs ---
\node[input] (mat)   {Таблица свойств пород\\\scriptsize{(Insar/data, гранит/базальт/...)}};
\node[input, right=of mat] (rules) {Правила сэмплирования\\\scriptsize{слои, источник, ГУ}};

% --- Stage 2: per-sample config ---
\node[proc, below=8mm of $(mat)!0.5!(rules)$] (cfg)
  {Конфигурация сэмпла $i$\\\scriptsize{$\{$слои, материалы, источник, ГУ$\}_i$}};

\draw[arr] (mat) -- (cfg);
\draw[arr] (rules) -- (cfg);

% --- Stage 3: COMSOL ---
\node[proc, below=of cfg] (comsol)
  {COMSOL \texttt{Untitled.mph} + Global Parameters\\
   \scriptsize{Parametric Sweep $\rightarrow$ N независимых решений}};
\draw[arr] (cfg) -- (comsol);

% --- Stage 4: raw export ---
\node[data, below=of comsol] (raw)
  {\textbf{Raw export на сэмпл} (\texttt{raw data/}):\\
   \scriptsize{\texttt{data\_displacement.csv},
   \texttt{data\_temperature.csv},
   \texttt{data\_stress\_\{1,2,3\}.csv},}\\
   \scriptsize{\texttt{data\_materials.csv}, \texttt{mesh.mphtxt}}};
\draw[arr] (comsol) -- (raw);

% --- Stage 5: preprocessing ---
\node[proc, below=of raw] (prep)
  {Препроцессинг \texttt{convert\_raw\_to\_graph.py}\\
   \scriptsize{разбор CSV $\to$ узлы, рёбра, признаки, нормализация}};
\draw[arr] (raw) -- (prep);

% --- Stage 6: per-sample npy ---
\node[data, below=of prep] (npy)
  {\textbf{Графовое представление сэмпла} (\texttt{data not raw/}):\\
   \scriptsize{\texttt{coords.npy} (N,3),\;
   \texttt{edge\_index.npy} (E,2),\;
   \texttt{edge\_attr.npy} (E,4),}\\
   \scriptsize{\texttt{node\_features.npy} (T,N,24),\;
   \texttt{times.npy} (T,),\;
   \texttt{norm\_mean/std.npy}}};
\draw[arr] (prep) -- (npy);

% --- Stage 7: master ---
\node[master, below=of npy] (master)
  {\textbf{Мастер-датасет}\\
   \scriptsize{\texttt{samples/sample\_XXXXX/} $\times$ N,\;
   \texttt{materials.csv},\; \texttt{meta.json},\; \texttt{splits.json}}};
\draw[arr] (npy) -- (master);

% --- Stage 8: four loaders ---
\node[loader, below=12mm of master, xshift=-48mm] (l1) {PINN-loader\\\scriptsize{collocation/BC/IC}};
\node[loader, below=12mm of master, xshift=-16mm] (l2) {FNO-loader\\\scriptsize{регулярная сетка}};
\node[loader, below=12mm of master, xshift=16mm]  (l3) {MGN-loader\\\scriptsize{узлы + рёбра, $t \to t{+}1$}};
\node[loader, below=12mm of master, xshift=48mm]  (l4) {Transformer-loader\\\scriptsize{input/query токены}};

\draw[arr] (master.south) -- ($(l1.north)+(0,2mm)$);
\draw[arr] (master.south) -- ($(l2.north)+(0,2mm)$);
\draw[arr] (master.south) -- ($(l3.north)+(0,2mm)$);
\draw[arr] (master.south) -- ($(l4.north)+(0,2mm)$);

% --- Stage 9: models ---
\node[model, below=6mm of l1] (m1) {PINN};
\node[model, below=6mm of l2] (m2) {FNO};
\node[model, below=6mm of l3] (m3) {MeshGraphNet};
\node[model, below=6mm of l4] (m4) {Transformer};

\draw[arrthin] (l1) -- (m1);
\draw[arrthin] (l2) -- (m2);
\draw[arrthin] (l3) -- (m3);
\draw[arrthin] (l4) -- (m4);

\end{tikzpicture}

\caption{Полный пайплайн: от случайной конфигурации среды до обучения
четырёх моделей. Этапы COMSOL и препроцессинга выполняются один раз
на сэмпл; адаптация под конкретную архитектуру~--- на лету в даталоадере.}
\label{fig:pipeline}
\end{figure}

\section{Текущее состояние и план}

\textbf{Текущее.} Реализован и проверен пилотный сэмпл~---
гомогенный гранитный куб, $1020$~узлов, $3001$~шаг по времени
($\Delta t = 0.01$~с), $24$~канала. Сырой экспорт~---
\texttt{Granite seismic simulation/raw data/}; графовое представление~---
\texttt{Granite seismic simulation/data not raw/}. Этот сэмпл играет роль
\emph{валидационного эталона} для будущих моделей.

\textbf{План.}
\begin{enumerate}[leftmargin=1.4em,itemsep=2pt]
\item Параметризация \texttt{Untitled.mph} через Global Parameters
      (материалы, слои, источник, ГУ).
\item Запуск \textit{Parametric Sweep} на 200--500 конфигураций
      (для дипломной работы достаточно).
\item Прогон \texttt{convert\_raw\_to\_graph.py} в режиме батча
      $\rightarrow$ заполнение \texttt{samples/sample\_*/}.
\item Реализация четырёх даталоадеров и базовое обучение по очереди:
      FNO~$\rightarrow$ MGN~$\rightarrow$ Transformer~$\rightarrow$ PINN.
\item Сравнительная таблица метрик на тестовой выборке
      (включая исходный гранитный сэмпл).
\end{enumerate}

\end{document}
